% Source excerpts; selections and clarifications documented in ../README.md.
% Medium: GPL v3; challenge: LPPL 1.3c.

\begin{defproblem}{T22-A01-08}%
\begin{onlyproblem}%
% \href{run:./ItemBankFigures/T22-A01-06.ggb}{（查看动态演示）}
已知，在$\triangle ABC$中，$\angle BAC=90^{\circ }$，$\angle ABC=45^{\circ }$，$D$为直线$BC$上一动点（不与点$B$，$C$重合），以$AD$为边作正方形$ADEF$，连接$CF$．

（1）如图1，当点$D$在线段$BC$上时，求证：$BC=CF+CD$．

（2）如图2，当点$D$在线段$BC$的延长线上时，其他条件不变，请直接写出$BC$，$CD$，$CF$三条线段之间的数量关系．

（3）如图3，当点$D$在线段$BC$的反向延长线上时，点$A$，$F$分别在直线$BC$的两侧，其他条件不变．

①请直接写出$BC$，$CD$，$CF$三条线段之间的数量关系；

②若正方形$ADEF$的边长为$2\sqrt 2$，对角线$AE$，$DF$相交于点$O$，求$OC$的长．


\begin{center}
\includegraphics[scale=0.8]{T22-A01-08-01.pdf}\qquad
\includegraphics[scale=0.8]{T22-A01-08-02.pdf}\qquad
\includegraphics[scale=0.8]{T22-A01-08-03.pdf}
\end{center}


% \begin{flushright}
% \includegraphics[scale=0.8]{T22-A01-02-01.pdf}\\
% \includegraphics[scale=0.8]{T22-A01-02-02.pdf}\\
% \includegraphics[scale=0.8]{T22-A01-02-03.pdf}
% \end{flushright}

\end{onlyproblem}%
\begin{onlysolution}%
（1）证明略．
提示：题目中有旋转结构，证明$\triangle ABD\cong \triangle ACF$(SAS)，
得到$BD=CF$，进而得证．

（2）$BC=CF-CD$．

（3）① $BC=CD-CF$；② $OC=2$．
	

\begin{center}
% \includegraphics[width=.95\textwidth]{T22-A01-06-04-Ans.jpg}
\end{center}
\end{onlysolution}%
\end{defproblem}

(1)如图a，四边形$ABCD,AB=CD$，$E,F$分别是$BC,AD$的中点，$EF$与$BA,CD$延长线交于$M,N$，求证：$\angle BME=\angle CNE$. \\
    (2)如图b，四边形$ADBC,AB=CD$，$E,F$分别是$BC,AD$的中点，连$EF$交$DC,AB$于$M,N$，判断$\triangle OMN$的形状，并证明. \\
    (3)如图c，$\triangle ABC,AC>AB$，$D$为$AC$上一点，且$AB=CD$，$E,F$分别是$BC,AD$的中点，连$EF$与$BA$延长线交于$G$.若$\angle EFC=60^\circ$，连$GD$，
    判断$\triangle AGD$的形状并证明.
\end{problem}
